\scanpage{52}
\margin{51}
- \mspage{73} - \hfill n$^{\mathrm{o}}$ 269

$\dim G<\dim F$, la restriction de $\Phi$ à $F\times G$ est dégénérée, et donc, comme $G\cap F^{\mathrm o}$ est nul, $F\cap G^{\mathrm o}$ est non nul. Si l'on avait alors $G+F^{\mathrm o}\supset G^{\mathrm o}$, on en déduirait, en prenant les orthogonaux et en remarquant que $F=F^{\mathrm{oo}}$ et que $G=G^{\mathrm{oo}}$, que $G^{\mathrm o}\cap F\subset G$, d'où $G^{\mathrm o}\cap F\subset G\cap F\subset G\cap F^{\mathrm o}=\{0\}$, ce qui est impossible. Il existe alors un élément $x$ de $G^{\mathrm o}$ tel que $x\notin G+F^{\mathrm o}$ ; comme $G\subset G^{\mathrm o}$ et $F\subset F^{\mathrm o}$, on peut ajouter à $x$ un vecteur de $G^{\mathrm o}\cap F$ sans modifier ces propriétés ; comme $G^{\mathrm o}\cap F$ est totalement isotrope (resp. totalement singulier) et $\ne\{0\}$, le lemme montre qu'on peut choisir $x$ isotrope (resp. singulier). Alors le sous espace $G'=G+Ax$ est de dimension $s+1$, et est totalement isotrope (resp. totalement singulier) ; de plus on a $G'\cap F^{\mathrm o}=\{0\}$ car, si $y=z+ax$ $(z\in G,a\in A)$ est dans $F^{\mathrm o}$, on a $a=0$ car sinon $x\in F^{\mathrm o}+G$ contrairement au choix de $x$, d'où $y=z\in G\cap F^{\mathrm o}=\{0\}$ et $y=0$. Par conséquent le sous espace $G'$ répond à la question.

COROLLAIRE - \uline{Si} $F$ \uline{est un sous-espace totalement isotrope} (resp. \uline{totalement singulier}) \uline{de dimension} $r$, \uline{il existe un sous espace totalement isotrope} (resp. \uline{totalement singulier}) $F'$ \uline{de dimension} $r$ \uline{tel que} $F\cap F'=\{0\}$ \uline{et que} $F+F'$ \uline{soit non isotrope}.

Il suffit de faire $G=\{0\}$ dans la prop. 2,b).

\uline{Remarque}. Dans les conditions du corollaire, $E$ est somme directe de $F+F'$ et de l'orthogonal de $F+F'$. On a donc une décomposition de Witt de $E$. D'après la prop.2,a), il existe des bases de $F$ et $F'$ telles que, dans la matrice (1) de $\Phi$, le bloc $\unsure{\mathcal U|\mathcal A}$ soit la matrice unité $\unsure{\underline{1}_r}$.
